\documentclass[11pt]{article}
\newcommand{\ListaTitulo}{Lista 05: Espaço amostral, eventos e contagem}
% Preâmbulo compartilhado: MATF14, Listas de Exercícios 2026
% Usado por ListaNN.tex e GabaritoNN.tex via % Preâmbulo compartilhado: MATF14, Listas de Exercícios 2026
% Usado por ListaNN.tex e GabaritoNN.tex via % Preâmbulo compartilhado: MATF14, Listas de Exercícios 2026
% Usado por ListaNN.tex e GabaritoNN.tex via \input{preamble}

\usepackage[utf8]{inputenc}
\usepackage[T1]{fontenc}
\usepackage[brazil]{babel}
\usepackage{fullpage}
\usepackage{amsmath,amssymb}
\usepackage{enumitem}
\usepackage{xcolor}
\usepackage{fancyhdr}
\usepackage{titlesec}
\usepackage{listings}
\usepackage{hyperref}
\usepackage{booktabs}
\usepackage{graphicx}
\usepackage{tikz}

\definecolor{matf14azul}{HTML}{0D3B54}
\definecolor{matf14dourado}{HTML}{C9971E}

\titleformat{\section}{\normalfont\Large\bfseries\color{matf14azul}}{\thesection}{1em}{}
\titleformat{\subsection}{\normalfont\large\bfseries\color{matf14azul}}{\thesubsection}{1em}{}

\providecommand{\ListaTitulo}{Lista de Exercícios}

\setlength{\headheight}{14pt}
\pagestyle{fancy}
\fancyhf{}
\lhead{\small MATF14: Estatística Econômica I}
\rhead{\small \ListaTitulo}
\cfoot{\thepage}
\renewcommand{\headrulewidth}{0.4pt}

\lstset{
  basicstyle=\ttfamily\small,
  breaklines=true,
  frame=single,
  columns=fullflexible,
  backgroundcolor=\color{gray!8},
  commentstyle=\color{matf14dourado},
  keywordstyle=\color{matf14azul}\bfseries,
  language=R
}

\newcounter{questaocounter}
\newenvironment{questao}[1]{%
  \stepcounter{questaocounter}%
  \bigskip\noindent\textbf{Questão #1.}\ %
}{\par}

\newcommand{\cabecalho}[2]{%
  \begin{center}
    {\Large\bfseries\color{matf14azul} #1}\\[0.3em]
    {\large #2}
  \end{center}
  \vspace{0.5em}
  \noindent\rule{\linewidth}{0.6pt}
  \vspace{0.5em}
}

\newenvironment{solucao}{%
  \par\smallskip\noindent\textit{\color{matf14dourado!80!black}Solução.}\ %
}{\hfill$\blacksquare$\par}


\usepackage[utf8]{inputenc}
\usepackage[T1]{fontenc}
\usepackage[brazil]{babel}
\usepackage{fullpage}
\usepackage{amsmath,amssymb}
\usepackage{enumitem}
\usepackage{xcolor}
\usepackage{fancyhdr}
\usepackage{titlesec}
\usepackage{listings}
\usepackage{hyperref}
\usepackage{booktabs}
\usepackage{graphicx}
\usepackage{tikz}

\definecolor{matf14azul}{HTML}{0D3B54}
\definecolor{matf14dourado}{HTML}{C9971E}

\titleformat{\section}{\normalfont\Large\bfseries\color{matf14azul}}{\thesection}{1em}{}
\titleformat{\subsection}{\normalfont\large\bfseries\color{matf14azul}}{\thesubsection}{1em}{}

\providecommand{\ListaTitulo}{Lista de Exercícios}

\setlength{\headheight}{14pt}
\pagestyle{fancy}
\fancyhf{}
\lhead{\small MATF14: Estatística Econômica I}
\rhead{\small \ListaTitulo}
\cfoot{\thepage}
\renewcommand{\headrulewidth}{0.4pt}

\lstset{
  basicstyle=\ttfamily\small,
  breaklines=true,
  frame=single,
  columns=fullflexible,
  backgroundcolor=\color{gray!8},
  commentstyle=\color{matf14dourado},
  keywordstyle=\color{matf14azul}\bfseries,
  language=R
}

\newcounter{questaocounter}
\newenvironment{questao}[1]{%
  \stepcounter{questaocounter}%
  \bigskip\noindent\textbf{Questão #1.}\ %
}{\par}

\newcommand{\cabecalho}[2]{%
  \begin{center}
    {\Large\bfseries\color{matf14azul} #1}\\[0.3em]
    {\large #2}
  \end{center}
  \vspace{0.5em}
  \noindent\rule{\linewidth}{0.6pt}
  \vspace{0.5em}
}

\newenvironment{solucao}{%
  \par\smallskip\noindent\textit{\color{matf14dourado!80!black}Solução.}\ %
}{\hfill$\blacksquare$\par}


\usepackage[utf8]{inputenc}
\usepackage[T1]{fontenc}
\usepackage[brazil]{babel}
\usepackage{fullpage}
\usepackage{amsmath,amssymb}
\usepackage{enumitem}
\usepackage{xcolor}
\usepackage{fancyhdr}
\usepackage{titlesec}
\usepackage{listings}
\usepackage{hyperref}
\usepackage{booktabs}
\usepackage{graphicx}
\usepackage{tikz}

\definecolor{matf14azul}{HTML}{0D3B54}
\definecolor{matf14dourado}{HTML}{C9971E}

\titleformat{\section}{\normalfont\Large\bfseries\color{matf14azul}}{\thesection}{1em}{}
\titleformat{\subsection}{\normalfont\large\bfseries\color{matf14azul}}{\thesubsection}{1em}{}

\providecommand{\ListaTitulo}{Lista de Exercícios}

\setlength{\headheight}{14pt}
\pagestyle{fancy}
\fancyhf{}
\lhead{\small MATF14: Estatística Econômica I}
\rhead{\small \ListaTitulo}
\cfoot{\thepage}
\renewcommand{\headrulewidth}{0.4pt}

\lstset{
  basicstyle=\ttfamily\small,
  breaklines=true,
  frame=single,
  columns=fullflexible,
  backgroundcolor=\color{gray!8},
  commentstyle=\color{matf14dourado},
  keywordstyle=\color{matf14azul}\bfseries,
  language=R
}

\newcounter{questaocounter}
\newenvironment{questao}[1]{%
  \stepcounter{questaocounter}%
  \bigskip\noindent\textbf{Questão #1.}\ %
}{\par}

\newcommand{\cabecalho}[2]{%
  \begin{center}
    {\Large\bfseries\color{matf14azul} #1}\\[0.3em]
    {\large #2}
  \end{center}
  \vspace{0.5em}
  \noindent\rule{\linewidth}{0.6pt}
  \vspace{0.5em}
}

\newenvironment{solucao}{%
  \par\smallskip\noindent\textit{\color{matf14dourado!80!black}Solução.}\ %
}{\hfill$\blacksquare$\par}


\begin{document}

\cabecalho{Lista de Exercícios 05}{Espaço amostral, eventos, axiomas de probabilidade e métodos de contagem (Cap. 2, Aulas 09--10)}

\noindent \textit{Justifique todas as respostas. As resoluções são discutidas em sala e nos horários de atendimento; não são publicadas neste material.}

\begin{questao}{1} \textbf{(Espaço amostral e eventos: carteira de dois ativos)}

Um investidor classifica o desempenho mensal de cada um de dois ativos, X e Y, como Alta (A),
Estável (E) ou Baixa (B).

\begin{enumerate}[label=(\alph*)]
  \item Escreva o espaço amostral $\Omega$ para o par (desempenho de X, desempenho de Y) em um
    mês. Quantos resultados possui?
  \item Escreva, como subconjunto de $\Omega$, o evento $C=$ "os dois ativos têm o mesmo
    desempenho no mês".
  \item Escreva o evento $D=$ "pelo menos um dos ativos está em Alta". $C$ e $D$ são mutuamente
    exclusivos? Justifique.
\end{enumerate}
\end{questao}

\begin{questao}{2} \textbf{(Axiomas e regra da adição: risco cambial e de juros)}

Uma empresa exportadora estima $P(\text{alta do dólar no trimestre})=0{,}45$,
$P(\text{alta de juros no trimestre})=0{,}30$, e $P(\text{ambos ocorrerem})=0{,}20$.

\begin{enumerate}[label=(\alph*)]
  \item Calcule a probabilidade de pelo menos um dos dois eventos (alta do dólar \emph{ou} alta de
    juros) ocorrer no trimestre.
  \item Calcule a probabilidade de \emph{nenhum} dos dois ocorrer.
  \item Um analista júnior calcula a resposta de (a) somando $0{,}45+0{,}30=0{,}75$, sem subtrair a
    interseção. Que evento está sendo contado duas vezes nesse cálculo errado, e por quê isso é um
    erro?
\end{enumerate}
\end{questao}

\begin{questao}{3} \textbf{(Contagem: formação de carteira de ações)}

Um gestor de fundos escolhe 4 ações para compor uma carteira, dentre 12 ações pré-aprovadas pelo
comitê de risco. A ordem da escolha não importa.

\begin{enumerate}[label=(\alph*)]
  \item Quantas carteiras diferentes de 4 ações são possíveis?
  \item Das 12 ações, 3 são do setor bancário. Quantas carteiras de 4 ações contêm \emph{exatamente}
    2 ações do setor bancário?
  \item Qual é a probabilidade de uma carteira de 4 ações, sorteada aleatoriamente entre as
    $\binom{12}{4}$ possíveis, conter exatamente 2 ações do setor bancário?
\end{enumerate}
\end{questao}

\begin{questao}{4} \textbf{(Princípio multiplicativo: senhas de acesso a um sistema bancário)}

Um sistema de internet banking exige uma senha numérica de 4 dígitos (cada dígito de 0 a 9, com
repetição permitida) seguida de uma pergunta de segurança escolhida entre 5 opções disponíveis.

\begin{enumerate}[label=(\alph*)]
  \item Quantas combinações possíveis de (senha, pergunta) existem?
  \item Se o banco proibir dígitos repetidos na senha (todos os 4 dígitos diferentes entre si),
    quantas senhas de 4 dígitos são possíveis? A proibição aumenta ou diminui a segurança contra
    tentativa aleatória de acesso? Justifique numericamente.
\end{enumerate}
\end{questao}

\begin{questao}{5} \textbf{(Dedução: a regra da adição geral)}

Sejam $A$ e $B$ dois eventos quaisquer (não necessariamente mutuamente exclusivos) em um espaço
amostral $\Omega$.

\begin{enumerate}[label=(\alph*)]
  \item Mostre que $A\cup B$ pode ser escrito como a união de três conjuntos mutuamente
    exclusivos: $A\cap B^c$, $A\cap B$, e $A^c\cap B$.
  \item Mostre também que $A = (A\cap B^c)\cup(A\cap B)$, uma união de dois conjuntos mutuamente
    exclusivos, e portanto $P(A) = P(A\cap B^c) + P(A\cap B)$.
  \item Usando os axiomas de aditividade para eventos mutuamente exclusivos e os itens (a)--(b),
    deduza a regra da adição geral: $P(A\cup B) = P(A) + P(B) - P(A\cap B)$. (Dica: escreva
    $P(A\cup B)$ usando o item (a), depois substitua $P(A\cap B^c)$ usando o item (b).)
\end{enumerate}
\end{questao}

\bigskip
\noindent\textit{A questão a seguir não tem resposta única e não é cobrada em prova no mesmo
formato, é para discussão em grupo ou por escrito, antes da próxima aula.}

\begin{questao}{6 (discussão)} \textbf{(Sem resposta única)}

Escolha um espaço amostral de um experimento econômico do seu interesse (por exemplo, resultado de
uma decisão de política monetária, desempenho trimestral de uma empresa, resultado de uma disputa
eleitoral). Defina $\Omega$ com um nível de detalhe que você julgue razoável, e depois argumente,
por escrito, por que uma definição mais grosseira (menos categorias) ou mais fina (mais categorias)
mudaria a análise que se pode fazer a partir dela.
\end{questao}


\bigskip
\section*{Exercícios complementares}

\begin{enumerate}[label=\textbf{C\arabic*.}, leftmargin=*]
\item Defina o espaço amostral de cada experimento:
\begin{enumerate}[label=(\alph*)]
  \item lançam-se dois dados e anota-se a configuração obtida;
  \item conta-se o número de peças defeituosas, em uma hora, de uma linha de produção;
  \item investigam-se famílias com três crianças e anota-se a configuração segundo o sexo;
  \item numa entrevista telefônica com cinco assinantes, pergunta-se se o proprietário tem ou não máquina de secar roupa;
  \item de um fichário com dois nomes de mulheres e dois de homens, sorteia-se uma ficha por vez até sair o último nome de mulher.
\end{enumerate}
\item Uma moeda é lançada três vezes. Seja $A_i$ o evento \emph{cara no $i$-ésimo lançamento}. Liste os elementos de:
\begin{enumerate}[label=(\alph*)]
  \item $A_1^c \cap A_3$
  \item $A_1^c \cup A_3$
  \item $(A_1^c \cap A_3^c)^c$
  \item $A_1 \cap (A_2 \cup A_3)$
\end{enumerate}
\item O espaço amostral é o intervalo $[0, 1]$. Sejam $A = \{x : 1/4 \leq x \leq 7/10\}$ e $B = \{x : 1/2 \leq x \leq 9/10\}$. Determine:
\begin{enumerate}[label=(\alph*)]
  \item $A^c$
  \item $A \cap B^c$
  \item $(A \cup B)^c$
  \item $A^c \cup B$
\end{enumerate}
\item Uma cidade tem 20 mil habitantes e três jornais, A, B e C. Uma pesquisa revela que 12 mil leem A, 8 mil leem B, 7 mil leem A e B, 6 mil leem C, 4.500 leem A e C, 1.000 leem B e C, e 500 leem os três. Sorteado um respondente, qual é a probabilidade de que ele leia
\begin{enumerate}[label=(\alph*)]
  \item pelo menos um jornal?
  \item somente um jornal? (Dica: faça o diagrama de Venn.)
\end{enumerate}
\item Estudantes de um curso de introdução à Economia, segundo o ano e se buscam especialização em Economia: 1 (2º, sim), 2 (4º, não), 3 (3º, sim), 4 (1º, não), 5 (1º, sim), 6 (2º, sim), 7 (2º, sim), 8 (3º, não), 9 (2º, sim), 10 (2º, não). Um estudante é escolhido ao acaso. Qual a probabilidade de ele ser
\begin{enumerate}[label=(\alph*)]
  \item do segundo ano?
  \item de especialização em Economia?
  \item do primeiro ou do segundo ano?
  \item de um ano diferente do primeiro?
\end{enumerate}
\item Um estudante responde ao acaso a duas questões: uma de verdadeiro/falso e outra de múltipla escolha com cinco alternativas. Qual a probabilidade de ele acertar
\begin{enumerate}[label=(\alph*)]
  \item a de verdadeiro/falso?
  \item a de múltipla escolha?
  \item ambas?
  \item nenhuma?
\end{enumerate}
\item De um grupo de 5 mulheres e 7 homens, quantos comitês diferentes com 2 mulheres e 3 homens podem ser formados? E se dois dos homens se recusam a trabalhar juntos?
\item Uma remessa de 150 celulares tem 40 defeituosos. Vinte são escolhidos ao acaso, sem reposição. Escreva a expressão, com combinações, da probabilidade de exatamente 9 defeituosos na amostra.

\end{enumerate}

\bigskip
\needspace{12\baselineskip}
\section*{Aprofundamento: novos tópicos e dados reais}
\noindent\textit{Exercícios sobre os tópicos acrescentados ao Capítulo 2 do livro e às aulas.}

\begin{enumerate}[label=\textbf{A\arabic*.}, leftmargin=*]

\item \textbf{(Interpretações de probabilidade.)} Diga qual interpretação (clássica, frequentista ou subjetiva) sustenta cada número:
\begin{enumerate}[label=(\alph*)]
  \item a chance de uma empresa específica ser sorteada para fiscalização entre 500 cadastradas é $1/500$;
  \item ``a probabilidade de inadimplência deste tipo de empréstimo é 3,2\%'', calculada com dez anos de contratos;
  \item ``há 70\% de chance de o Copom cortar a Selic na próxima reunião'', segundo um analista;
  \item um analista aceita comprar ou vender, por $p$ reais, bilhetes que pagam R\$ 1 se o evento ocorrer. Ele atribui
    $P(\text{Selic sobe}) = 0{,}55$ e $P(\text{Selic não sobe}) = 0{,}50$. Mostre como um investidor pode ganhar dinheiro com certeza.
\end{enumerate}

\item \textbf{(Axiomas de Kolmogorov.)} Usando apenas os axiomas e as propriedades vistas em aula:
\begin{enumerate}[label=(\alph*)]
  \item mostre que $P(A \cap B^c) = P(A) - P(A \cap B)$;
  \item mostre que $P(A \cap B) \geq P(A) + P(B) - 1$ (desigualdade de Bonferroni);
  \item em uma pesquisa, 65\% das empresas exportam e 55\% importam. Qual a menor e qual a maior proporção possível de
    empresas que fazem as duas coisas?
\end{enumerate}

\item \textbf{(Os quatro tipos de contagem.)} Para cada situação, diga se a ordem importa e se há repetição, e conte:
\begin{enumerate}[label=(\alph*)]
  \item senhas de 6 caracteres formadas com 26 letras e 10 algarismos, podendo repetir;
  \item ordens possíveis de apresentação de 7 propostas em uma licitação;
  \item comissões de 3 pessoas escolhidas entre 12 servidores;
  \item escolhas de presidente, vice e tesoureiro entre os mesmos 12 servidores;
  \item formas de distribuir 8 bolsas de estudo idênticas entre 4 cursos (algum curso pode ficar sem bolsa).
\end{enumerate}

\item \textbf{(Auditoria.)} Um lote tem 50 contratos, dos quais 5 são irregulares. O auditor sorteia 6 contratos.
\begin{enumerate}[label=(\alph*)]
  \item Sem reposição, qual a probabilidade de nenhum contrato irregular na amostra? E de pelo menos um?
  \item Refaça (a) supondo sorteio \emph{com} reposição e compare. Por que os valores são próximos?
\end{enumerate}

\item \textbf{(Quando contar não basta.)} Um colega afirma: "ao lançar três moedas, o número de caras pode ser 0, 1, 2 ou 3;
logo $P(\text{duas caras}) = 1/4$``. Aponte o erro e calcule a probabilidade correta. Faça o mesmo para ''a soma de dois dados
vai de 2 a 12, logo $P(\text{soma } 7) = 1/11$".

\end{enumerate}

\end{document}
